\subsection{模型建立}
我们在问题五中沿用问题二的状态方程，并在污染强度、入侵初值和竞争强度上做情景扩展。由于问题一的基线轨迹对应 $u=0$，污染情景下的上游输入需要再次由问题一积分得到，以便使资源退化、鱼群变化与污染设定保持一致。具体而言，先以 $u>0$ 重新计算 $E_{\text{basal}}$ 和 $E_{\text{carp}}$，再将其输入问题二模块。

考虑到污染和入侵会诱发幅度较大的波动，我们统一比较积分后期最后 200 个离散采样点上的平均指标，以减弱瞬时相位的影响。设第 $i$ 个情景在该尾窗上的四个平均指标分别为 $\bar M_i,\bar D_i,\bar A_i,\bar V_i$，其中食源、江豚和长江鲟为效益型指标，入侵压力为成本型指标。先对三种情景构造正向化指标矩阵，再分别计算熵权和 CRITIC 权重，并取两者平均作为组合权重。记组合权重为 $w=(w_1,w_2,w_3,w_4)$，则综合得分定义为
\begin{equation}
\mathcal{H}_i=\sum_{j=1}^{4}w_j z_{ij},\qquad \sum_{j=1}^{4}w_j=1,\quad w_j\ge 0.
\end{equation}
在当前数据下，组合权重为 $(0.232,0.289,0.298,0.181)$，分别对应食源、江豚、长江鲟和入侵压力。该指标用于比较三种情景的相对状态，不对应外部生态健康的绝对量值。

